\begin{frame}{프로젝트 구조}
  \begin{itemize}
    \item \kb{데이터} — Kaggle, UCI 외에 \kb{Hugging Face Datasets}도 후보
      (Ch11~13 시퀀스/LLM 프로젝트는 텍스트를 몇 줄로 불러옴)
    \item \kb{적용} — Ch09~15 중 \textbf{최소 하나의 딥러닝 기법}
      (CNN, RNN/Transformer, 프롬프팅 응용, 또는 EM/VAE 계열 생성모델)
     을 실제 데이터에 적용
    \item \kb{필수: 수식적 정당화 최소 1개} —
      (1) CNN 구조면 Ch10.1 파라미터 수/연산량 공식으로 설계 정당화,
      (2) 정규화(dropout 등)면 Ch9.3 효과 설명,
      (3) 생성모델이면 Ch15 ELBO나 min-max 목적함수를 실제 구현에 연결
    \item \kb{GBDT 비교(필수)} — 같은 데이터에 Ch7.3 GBDT(또는
      랜덤포레스트)를 \textbf{베이스라인}으로 돌려, 신경망이 이보다
      나은지 확인. 발표 ``왜 DL'' 근거의 출발점
    \item train/validation/test 분리(Ch06.3)를 8.1과 동일하게 준수
  \end{itemize}
\end{frame}
